\documentclass[11pt]{article}
\usepackage{fontspec}
\setmainfont{Times New Roman}
\setsansfont{Arial}
\setmonofont{Consolas}
\usepackage{amsmath,amssymb}
\usepackage{geometry}
\usepackage{microtype}
\geometry{a4paper,margin=3cm}
\setlength{\parindent}{1.25em}
\setlength{\parskip}{0pt}
\emergencystretch=4em
\allowdisplaybreaks
\newcommand{\eps}{\varepsilon}
\newcommand{\pair}[2]{\left(#1\,\middle|\,#2\right)}

\begin{document}

\begin{center}
{\Large\bfseries 3. До інваріантної теорії форм від $n$ змінних\footnote{Доповідь, виголошена на Зальцбурзькому з'їзді природознавців 1909 року.}}\par
\vspace{1.2em}
\emph{Jahresbericht d. DMV} 19 (1910), с. 101--104
\end{center}

\vspace{1em}
У проективній інваріантній теорії форм від $n$ змінних головну проблему розв'язано: скінченність системи форм доведено. Проте низку дальших питань ще не розглянуто, а саме тих, що стосуються методів дослідження зв'язку між формами. Результати, які я одержала щодо таких питань, я хочу коротко подати, але для ясності спершу окреслю самі постановки питань у тернарній області, де вони вже відомі.

Я маю на увазі насамперед дві фундаментальні теореми символічного методу: теорему про те, що всі інваріантні утворення можна подати символічно, і другу теорему, за якою всі відношення між інваріантами одержуються послідовним застосуванням невеликої кількості символічних тотожностей; отже, символічний метод дає всі відношення, не виходячи з області інваріантів.\footnote{Праця: \emph{Methoden zur Theorie der ternären Formen}. II. \S\ 6. (Leipzig, Teubner, 1889.)} На цьому грунтуються процеси утворення форм, символічний процес згортання і несимволічні процеси диференціювання, теорія розкладів у ряди тощо. До цього ж я хотіла б додати ще одну теорему, доведену в моїй дисертації,\footnote{\emph{Über die Bildung des Formensystems der ternären biquadratischen Form}. \S\ 1. (\emph{Crelle's Journal}, Bd. 134, 1908.)} а саме: всі згортання можна породити послідовним застосуванням двох можливих ``когредієнтних'' згортань; під цим для символічного добутку $s_x t_x u_\sigma u_\tau$ я розумію два згортання $(stu)$ і $(\sigma\tau x)$. На цій теоремі можна побудувати теорію редуцентів і редукції систем форм, аналогічно до бінарної області, як я показала у своїй дисертації.

Якщо тепер перейти до $n$-арної області, то вже перша теорема символічного методу справджується лише умовно. Як відомо, вже для кватернарних форм, крім точкових координат $x$ і координат площин $u$, виникають координати прямих $p_{ik}$, тобто мінори з двох рядків точкових координат. Аналогічно в $n$-арній області маємо рядки змінних $p_\rho$, окремі елементи яких є мінорами з $\rho$ рядків $x$, і символьні рядки $S^\rho$, елементи яких поводяться як мінори з $\rho$ рядків $s$, когредієнтних до $u$. Символічне подання через детермінанти, як відомо, чинне лише тоді, коли символьні рядки $S^\rho$ розкладено на рядки $s$ і ці рядки розподілено між різними детермінантами; тоді реальний зміст мають тільки такі агрегати детермінантів, які залежать від $S^\rho$. Але неможливо наперед оглянути, які саме агрегати детермінантів це забезпечують; через це втрачаються всі інші теореми, згадані для тернарної області.

Тепер я хочу сформулювати простий принцип, який дає символічне подання, явне щодо рядків $S^\rho$, і тим самим відновлює також решту теорем. Це застосування відомої теореми про відповідні матриці: ``Кожному рядку змінних $p_\rho$ взаємно однозначно відповідає інший рядок $q_{n-\rho}$, окремі елементи якого є мінорами з $n-\rho$ рядків $u$, пов'язаних із $\rho$ рядками $x$ рівняннями $(u\mid x)=0$.'' Відповідно кожному символьному рядку $S^\rho$ зіставлено інший рядок $S_{n-\rho}$, який поводиться так, ніби він складений із $n-\rho$ рядків, когредієнтних до $x$.

Цю теорему можна сформулювати ще й у другому, зручному для застосувань вигляді: ``Кожному матричному добутку $(v^{(1)}\cdots v^{(\rho)}\mid y^{(1)}\cdots y^{(\rho)})$,\footnote{Тобто сумі за добутками відповідних детермінантів, утворених із матриці $\rho$ рядків $v$ і матриці $\rho$ рядків $y$.} де рядки другої матриці контрагредієнтні до рядків першої, відповідає детермінант; і навпаки, кожен детермінант можна перетворити на матричний добуток із наперед заданою кількістю рядків $\rho$.'' Отже, детермінант і матричний добуток постають як цілком еквівалентні, а перша теорема символічного методу набуває форми: ``Усі інваріантні утворення можна символічно подати матричними добутками.'' З двох символьних рядків $S^\sigma$, $T^\tau$ за допомогою рядків змінних тепер дуже легко утворити матричний добуток, що явно містить ці символьні рядки і тому подає інваріантне утворення форми $(S^\sigma\mid p_\sigma)(T^\tau\mid p_\tau)$, а саме:
\begin{equation}
 \pair{q_{n-\sigma-\lambda}S^\sigma}{T_{n-\tau}p_{\tau-\lambda}},
 \qquad (\lambda=1,2,\ldots,\tau\ \text{або}\ n-\sigma).
\tag{1}
\end{equation}

Важливішим є обернене твердження: всі інваріантні утворення можна явно подати матричними добутками, отже наведений процес є розширенням бінарного й тернарного процесу згортання. Щоб довести це обернення, треба перенести другу теорему символічного методу на $n$-арну область, тобто виписати всю сукупність незалежних нерозкладних тотожностей, у яких усі рядки $S^\rho$, $p_\rho$ розглядаються як нерозкладні. Ці тотожності одержуються з відомих тотожностей, що містять лише рядки $x$ і $u$,\footnote{E. Pascal у \emph{Memorie d. R. Acc. d. Lincei}. (4) 5 (1888), p. 375.} якщо замінити теорему про добуток для детермінантів системою теорем про добуток для матриць і стягувати різні матричні добутки в один саме за допомогою цих теорем про добуток. Так отримуємо дві дуалістичні формули, з яких наведу тільки одну:
\begin{equation}
\small
\begin{aligned}
 \pair{S_1^{\rho_1}S_2^{\rho_2}\cdots S_k^{\rho_k}}{x p_{\rho-1}}
 &=\sum_{i=1}^{k}\eps\,\pair{S_1^{\rho_1}\cdots S_{i-1}^{\rho_{i-1}}S_{i+1}^{\rho_{i+1}}\cdots S_k^{\rho_k}q_{n-\rho_i+1}}{S_{n-\rho_i}x},\\
 \rho&=\sum_{i=1}^{k}\rho_i<n,\qquad \eps=\pm1.
\end{aligned}
\tag{2}
\end{equation}

Єдина спеціалізація, що міститься в цих формулах, а саме поява рядка $x$, є неминучою; для цілком загальних символьних рядків і рядків змінних ми приходимо до ``тотожності розкладання'', тобто до тотожності, в якій рядок $p_\rho$ розкладається на окремі рядки $x$. Найпростіший окремий випадок тотожності розкладання вже використав Клебш у своїй ``фундаментальній задачі інваріантної теорії'' для виведення елементарних коваріантів у теорії розкладів у ряди. За допомогою цих тотожностей, які здійснюють перехід від неявних детермінантних виразів до матричного добутку, справді можна довести, що матричне подання дає бажане явне символічне подання.

Для процесу згортання, однак, справджується теорема, цілком аналогічна до тернарного випадку; на ній так само аналогічно можна побудувати теореми редукції. Якщо в (1) назвати число $\lambda$ дефектом згортання, а згортання, для яких $\sigma=\tau$, назвати когредієнтними, то всі згортання можна породити послідовним застосуванням когредієнтних згортань дефекту 1. Доведення цієї теореми істотно спирається на те, що найзагальніші форми можна замінити ``нормальними формами'', тобто формами, які тотожно зникають під дією всіх інваріантних диференціальних операцій. У кватернарній області цей останній факт довів Мертенс;\footnote{\emph{Über invariante Gebilde quaternärer Formen}. \emph{Wiener Berichte}. Bd. 98, 1889.} наше знання найзагальніших диференціальних операцій -- які, як відомо, виникають із процесів згортання через заміну символьних рядків диференціальними символами -- дозволяє, застосувавши тотожність розкладання, майже дослівно перенести доведення Мертенса на $n$-арну область.

\end{document}
